\section{A multi-core CPU baseline: the OpenMP version}\label{sec:omp}

Comparing one CPU core with a whole GPU is an unfair contest, and a reader can fairly ask whether the GPU wins
only because it is parallel. We therefore parallelised the sequential program (\texttt{seq/rasp\_seq.cpp}) with OpenMP, changing no algorithm, so
that all CPU cores can be compared with the GPU. The sequential, OpenMP and CUDA versions read the same inputs and
print the same phase breakdown. All numbers here come from one laptop CPU (an Intel Core i5-1235U: two
performance cores with two hardware threads each and eight efficiency cores, 12 hardware threads, under WSL2
Ubuntu, g++ 15 with \texttt{-O2 -march=native}), with the sequential baseline measured on the same machine.
Threads are pinned to separate cores (\texttt{OMP\_PROC\_BIND=close}, \texttt{OMP\_PLACES=cores}).

\subsection{What was parallelised}
\begin{table}[t]
\centering\footnotesize
\caption{Kernels of \texttt{seq/rasp\_omp.cpp} and how each is parallelised.}\label{tab:omp-kernels}
\begin{tabular}{@{}p{2.1cm}p{5.8cm}@{}}
\toprule
kernel & OpenMP \\
\midrule
SpMV CSR, SELL & \texttt{parallel for} over rows (CSR) or slices (SELL); loop schedule chosen at run time (\texttt{--schedule}) \\
SpMV ELL & one contiguous block of rows per thread, each running the original column-by-column loops (no fork/join per column) \\
\texttt{dot} & \texttt{reduction(+)}; the summation order changes, so CG may differ from the sequential run in rounding only \\
\texttt{axpy}, Jacobi, $p$-update, CG diagonal & \texttt{parallel for schedule(static)} \\
SpGEMM flop estimate & \texttt{parallel for} with a \texttt{reduction(+)} \\
SpGEMM symbolic pass & per-thread marker array; each row stores only its count, the prefix sum is sequential ($O(\text{rows})$) \\
SpGEMM numeric pass & per-thread accumulator, marker and list of touched columns; row $i$ writes only its own slice $[rp_i, rp_{i+1})$ of $C$, so there are no conflicts \\
\bottomrule
\end{tabular}
\end{table}

Rows have very different lengths, which makes the loop schedule matter for SpMV on irregular matrices: this is the
CPU version of the load-balancing problem ArgCSR addresses on the GPU. The SpGEMM loops use
\texttt{schedule(dynamic, 256)} for the same reason. File parsing and format conversion stay sequential, because
the comparisons do not time them.

\subsection{Correctness}
\texttt{rasp\_omp --verify} keeps the original sequential kernels and compares against them on the same inputs, at
1, 2 and all threads. SpMV in every format must agree to a relative error of $10^{-12}$; CG must converge to the
same tolerance within two iterations of the sequential count; and every AMG coarse matrix must have exactly the
sequential \texttt{row\_ptr} and \texttt{col}, with values equal to $10^{-12}$. All checks pass (in practice SpMV
CSR matches exactly, the other formats differ by about $10^{-16}$, and the CG and AMG results are identical).

\subsection{Scaling}
\begin{figure}[t]
\centering\rfig[0.62\linewidth]{omp_scaling}
\caption{Speedup over the sequential run against thread count (fastest of two sweeps). The dashed line is the
ratio of all-thread to single-thread streaming bandwidth (\stat{omp-bw-1}\,GB/s to \stat{omp-bw-all}\,GB/s,
\stat{omp-bw-ceiling}\x), the expected ceiling for memory-bound kernels.}\label{fig:omp-scaling}
\end{figure}

\begin{table}[t]
\centering\small
\caption{Speedup over the sequential version on the same machine (fastest of two sweeps). SpMV is the geometric
mean over six matrices; CG solves a $128^3$ Laplacian; AMG is the SpGEMM time of a $160^3$ Laplace hierarchy.}\label{tab:omp-scaling}
\stattable{omp-scaling}
\end{table}

\Cref{tab:omp-scaling} and \cref{fig:omp-scaling} show the result. Speedups stop at about \stat{omp-spmv-best-speedup}\x{} (SpMV, \stat{omp-spmv-best-threads}
threads), \stat{omp-cg-best-speedup}\x{} (CG, \stat{omp-cg-best-threads} threads) and \stat{omp-amg-best-speedup}\x{} (AMG SpGEMM,
\stat{omp-amg-best-threads} threads), against a hardware limit of 12. Within one sweep, where no baseline drift enters, the
best gains over the same code on one thread are \stat{omp-spmv-best-speedup-vs1}\x, \stat{omp-cg-best-speedup-vs1}\x{} and
\stat{omp-amg-best-speedup-vs1}\x. The OpenMP version on one thread runs as fast as the sequential version (back-to-back runs agree
within measurement noise), so the OpenMP overhead itself is negligible.

\textbf{Memory bandwidth is the ceiling for SpMV and CG.} These kernels do two flops per twelve bytes read, so they
are limited by memory bandwidth, and a CPU saturates it with few threads: one thread streams \stat{omp-bw-1}\,GB/s,
all \stat{omp-bw-threads} threads \stat{omp-bw-all}\,GB/s. CG, whose time is
dominated by SpMV and by streaming vector updates, flattens at about this ratio. SpMV reaches slightly above it on some matrices, plausibly because part of their
working set stays in cache. For comparison the RTX~4050 reads \stat{pf-device-read}\,GB/s: the GPU's advantage on these kernels
is mainly bandwidth, which no number of CPU cores supplies.

\textbf{The loop schedule is not resolved at this noise level.} \Cref{tab:omp-sched} compares the loop schedules at the highest thread
count. We would expect \texttt{dynamic} to help on matrices with a few very long rows (\texttt{rajat23}), as the
GPU's row grouping does; at this thread count the three schedules are dominated by run-to-run noise, so we draw no
conclusion about their ranking.

\begin{table}[t]
\centering\small
\caption{SpMV CSR with each loop schedule at the highest thread count (geometric mean over six matrices).}\label{tab:omp-sched}
\stattable{omp-schedule}
\end{table}

\textbf{SpGEMM is held back by sequential system time, not by memory bandwidth.} We expected the AMG set-up to scale
further than SpMV, since it is hash and accumulator work. It does not. In the numeric pass, \texttt{std::vector::resize}
zero-fills the output arrays on a single thread, and the kernel page faults that this causes are charged as system time:
\stat{omp-amg-numeric-sys-s-t1}\,s (\stat{omp-amg-numeric-sys-pct-t1}\% of the pass) on one thread
and \stat{omp-amg-numeric-sys-s-best}\,s (\stat{omp-amg-numeric-sys-pct-best}\% of the pass) at the best thread count. That time does not shrink with the thread count, so
Amdahl's law caps the speedup. The prefix sum of the symbolic pass is sequential too, but it is negligible next to this.

\textbf{Limits of this measurement.} The i5-1235U has two kinds of core and the sequential run may land on a
performance core with a high turbo clock, which makes parallel efficiency look lower than it is. The 12-thread
results use hyperthread siblings and efficiency cores and varied much more between the two sweeps than the
1--8 thread results, so we do not rely on them; sequential-baseline drift between runs is about $\pm$10\%.
WSL2 also inflates system time. The same sweep runs on the Kaggle and RTX~4050 machines, so every
comparison is against a sequential and a GPU run on the same hardware.

\subsection{OpenMP against the GPU}
Even with every core busy, the CPU version is limited to \stat{omp-cg-best-speedup}--\stat{omp-spmv-best-speedup}\x{} over one core on its memory-bound kernels,
because its bandwidth rises only \stat{omp-bw-ceiling}\x{} with all threads, while the GPU has several times the bandwidth.
\Cref{tab:same-machine} compares the GPU with the CPU on the same machines.
